\documentclass[12pt,a4paper]{article}
\usepackage[french]{babel}
\usepackage[utf8]{inputenc}
\usepackage[T1]{fontenc}
\usepackage{geometry}
\usepackage{graphicx}
\usepackage{booktabs}
\usepackage{array}
\usepackage{xcolor}
\usepackage{hyperref}
\usepackage{fancyhdr}
\usepackage{titlesec}
\usepackage{amsmath}
\usepackage{listings}
\usepackage{caption}
\usepackage{subcaption}
\usepackage{float}
\usepackage{multirow}

\geometry{top=2.5cm, bottom=2.5cm, left=2.5cm, right=2.5cm}

% Couleurs
\definecolor{maingreen}{RGB}{34,197,94}
\definecolor{darkbg}{RGB}{15,23,42}
\definecolor{codegray}{RGB}{245,245,245}

% Titres
\titleformat{\section}{\large\bfseries\color{darkbg}}{}{0em}{\thesection\quad}[\titlerule]
\titleformat{\subsection}{\normalsize\bfseries}{}{0em}{\thesubsection\quad}

% En-tête / Pied de page
\pagestyle{fancy}
\fancyhf{}
\fancyhead[L]{\small\textit{Détection de Maladies des Plantes}}
\fancyhead[R]{\small\textit{IGL4 — 2025/2026}}
\fancyfoot[C]{\thepage}
\renewcommand{\headrulewidth}{0.4pt}

% Listings
\lstset{
  backgroundcolor=\color{codegray},
  basicstyle=\ttfamily\footnotesize,
  breaklines=true,
  frame=single,
  language=Python,
  keywordstyle=\color{blue},
  commentstyle=\color{gray},
  showstringspaces=false
}

\hypersetup{colorlinks=true, linkcolor=darkbg, urlcolor=maingreen}

% ============================================================
\begin{document}

% ============================================================
% PAGE DE GARDE
% ============================================================
\begin{titlepage}
\centering
\vspace*{1cm}
{\large\textbf{Université Privée de Tunis}\\
Institut des Hautes Études Technologiques\\[0.3cm]
Section IGL4 — Année universitaire 2025/2026}

\vspace{2cm}
\rule{\linewidth}{2pt}\\[0.4cm]
{\Huge\bfseries\color{darkbg}
Détection et Classification\\[0.2cm]
Automatique de Maladies\\[0.2cm]
des Plantes}\\[0.4cm]
\rule{\linewidth}{2pt}

\vspace{1.5cm}
{\large\textit{Projet de Traitement d'Images \& Intelligence Artificielle}}

\vspace{2cm}
\begin{tabular}{rl}
\textbf{Auteurs :} & Bahaeddine ELLOUZE \\
                   & Omar CHOKRI \\[0.3cm]
\textbf{Section :} & IGL4 \\[0.3cm]
\textbf{Encadrant :} & \textit{[Nom de l'encadrant]} \\[0.3cm]
\textbf{Date :} & Avril 2026 \\
\end{tabular}

\vfill
{\small\textit{Rapport remis dans le cadre du module Traitement d'Images}}
\end{titlepage}

% ============================================================
% TABLE DES MATIÈRES
% ============================================================
\tableofcontents
\newpage

% ============================================================
\section{Introduction}
% ============================================================

Les maladies des plantes représentent une menace majeure pour la sécurité alimentaire mondiale, 
causant des pertes agricoles estimées entre 20\% et 40\% de la production annuelle 
(FAO, 2019). La détection précoce est essentielle pour limiter les dommages, 
mais elle requiert une expertise agronomique difficile à diffuser à grande échelle.

Ce projet propose un \textbf{système intelligent de détection et classification automatique} 
de maladies foliaires à partir d'images, combinant vision par ordinateur classique, 
Machine Learning et Deep Learning. Il est réalisé par \textbf{Bahaeddine Ellouze} et 
\textbf{Omar Chokri} dans le cadre du module de traitement d'images (IGL4).

\subsection{Objectifs}
\begin{itemize}
  \item Implémenter un pipeline complet de bout en bout (prétraitement $\rightarrow$ prédiction)
  \item Comparer les approches ML classique (RF, SVM) et DL (MobileNetV2)
  \item Fournir une interface utilisateur interactive et une explicabilité visuelle (Grad-CAM)
\end{itemize}

\subsection{Dataset utilisé}
Nous utilisons le \textbf{PlantVillage Dataset} (Kaggle), en sélectionnant 
\textbf{12 classes} sur 3 espèces végétales (binôme $\geq$ 8 classes requises) :

\begin{table}[H]
\centering
\small
\begin{tabular}{lll}
\toprule
\textbf{Espèce} & \textbf{Classes} & \textbf{Images} \\
\midrule
Maïs (Corn) & Cercospora, Common rust, Northern Blight, Saine & 4 188 \\
Pomme de terre & Early blight, Late blight, Saine & 2 152 \\
Tomate & Early blight, Late blight, Leaf Mold, Septoria, Saine & 7 223 \\
\midrule
\textbf{Total} & \textbf{12 classes} & \textbf{13 563 images} \\
\bottomrule
\end{tabular}
\caption{Distribution du dataset PlantVillage utilisé}
\end{table}

% ============================================================
\section{Architecture Globale du Pipeline}
% ============================================================

Notre pipeline suit une architecture modulaire en 6 étapes séquentielles :

\begin{figure}[H]
\centering
% Insérer ici un schéma du pipeline (ex: class_balance.png)
\fbox{\parbox{0.85\textwidth}{\centering\vspace{1cm}
\textit{[Figure : Schéma du pipeline complet — de l'image brute à la prédiction]}
\vspace{1cm}}}
\caption{Pipeline complet de détection et classification}
\end{figure}

\noindent Le code source est organisé en modules Python dans le répertoire \texttt{src/} :
\texttt{data/} (chargement), \texttt{vision/} (segmentation), \texttt{features/} (extraction), 
\texttt{utils/} (benchmark, conseils). L'entraînement est lancé via \texttt{train\_ml.py} 
et \texttt{train\_dl.py}.

% ============================================================
\section{Prétraitement des Images}
% ============================================================

\subsection{Redimensionnement (224 $\times$ 224 pixels)}
Toutes les images sont redimensionnées à $224 \times 224$ pixels avec l'interpolation 
\texttt{INTER\_AREA} d'OpenCV, optimale pour la réduction d'image. Cette taille est imposée 
par MobileNetV2 (standard ImageNet) et assure la cohérence entre le pipeline ML et DL.

\subsection{Filtrage Gaussien (noyau 5$\times$5)}
Un filtre Gaussien ($k=5$, $\sigma_{\text{auto}}=1.25$) est appliqué pour supprimer le 
bruit de capteur photographique et les artéfacts de compression JPEG. Le noyau $5\times5$ 
préserve les contours des taches foliaires ($>$3 px de rayon) tout en lissant le bruit 
de capteur ($\sim$1--2 px).

\begin{figure}[H]
\centering
\fbox{\parbox{0.9\textwidth}{\centering\vspace{1.5cm}
\textit{[Figure : Avant / Après filtrage Gaussien — zoom sur un patch 64$\times$64]}
\vspace{1.5cm}}}
\caption{Impact du filtre Gaussien 5$\times$5 sur une image de Tomato Late blight}
\end{figure}

\subsection{Conversion RGB $\rightarrow$ HSV}
L'espace HSV (Teinte / Saturation / Valeur) sépare la teinte de la luminosité, 
rendant la segmentation robuste aux variations d'éclairage. La teinte H reste 
stable ($H \approx 60$--$90°$ pour le vert) quelle que soit l'intensité lumineuse.

\subsection{Analyse des Histogrammes}
Nous comparons les histogrammes du canal H (Teinte) de 4 classes, révélant une 
\textbf{signature spectrale distincte} par maladie : feuilles saines ($H \approx 60$--$90°$), 
rouille ($H \approx 20$--$30°$), nécrose ($H \approx 5$--$22°$).

\subsection{Partitionnement Stratifié 70/15/15}
Le dataset est divisé en Train (9\,490), Validation (2\,029) et Test (2\,044) 
avec une graine fixe (\texttt{seed=42}) et un nettoyage préalable de \texttt{data/processed/} 
pour éviter toute fuite de données (data leakage).

% ============================================================
\section{Segmentation et Détection de Contours}
% ============================================================

Nous avons implémenté et comparé \textbf{5 méthodes} de segmentation/détection.

\subsection{Détection de Contours — Sobel et Canny}

\textbf{Sobel} calcule le gradient d'intensité : $G = \sqrt{G_x^2 + G_y^2}$ 
avec un noyau $3\times3$.

\textbf{Canny} améliore Sobel par non-maxima suppression et double seuillage par 
hystérésis ($T_{\text{low}}=50$, $T_{\text{high}}=150$, ratio 3:1). Il produit 
des contours fins et précis.

\begin{figure}[H]
\centering
\fbox{\parbox{0.9\textwidth}{\centering\vspace{1.5cm}
\textit{[Figure : Comparaison Sobel / Canny / HSV sur Tomato Late blight]}
\vspace{1.5cm}}}
\caption{Détection de zones anormales — Sobel vs Canny vs référence HSV}
\end{figure}

\subsection{Seuillage d'Otsu}
L'algorithme d'Otsu calcule automatiquement le seuil optimal $T^*$ qui maximise 
la variance inter-classes fond/feuille. Il ne nécessite aucun paramètre manuel.

\subsection{Clustering K-Means ($k=3$)}
K-Means regroupe les pixels en 3 clusters (fond, feuille, taches) selon leur 
distance Euclidienne dans l'espace RGB. L'initialisation K-Means++ garantit 
la stabilité de la convergence.

\subsection{Double Masque HSV — Méthode Principale}
Notre pipeline ML utilise un double masque HSV calibré biologiquement :
\begin{itemize}
  \item \textbf{Masque vert} : $H \in [22°, 95°]$, $S > 30$, $V > 30$ — feuilles saines
  \item \textbf{Masque brun} : $H \in [5°, 22°]$, $S > 50$, $V < 200$ — zones nécrotiques
\end{itemize}
Suivi d'opérations morphologiques OPEN + CLOSE (noyau elliptique $7\times7$).

\begin{table}[H]
\centering\small
\begin{tabular}{lccccc}
\toprule
\textbf{Méthode} & \textbf{Isolation feuille} & \textbf{Détection taches} & \textbf{Vitesse} & \textbf{Paramètres} \\
\midrule
Sobel & Non & Partielle & Très rapide & ksize \\
Canny & Non & Bonne & Rapide & $T_l$, $T_h$ \\
Otsu & Oui (fond clair) & Limitée & Très rapide & Aucun \\
K-Means & Oui & Partielle & Lente & $k$ \\
\textbf{HSV masque} $\star$ & \textbf{Oui} & \textbf{Très bonne} & \textbf{Rapide} & \textbf{Seuils} \\
\bottomrule
\end{tabular}
\caption{Comparaison subjective des 5 méthodes de segmentation}
\end{table}

% ============================================================
\section{Extraction de Caractéristiques}
% ============================================================

Chaque image segmentée est transformée en un \textbf{vecteur de 566 dimensions} :

\begin{table}[H]
\centering\small
\begin{tabular}{llcp{5cm}}
\toprule
\textbf{Catégorie} & \textbf{Méthode} & \textbf{Dims} & \textbf{Capture quoi ?} \\
\midrule
Couleur RGB & Histogramme 1D (8 bins $\times$ 3) & 24 & Distribution des canaux R, G, B \\
Couleur HSV & Histogramme 3D ($8^3$ bins) & 512 & Teinte/saturation invariante éclairage \\
Texture & GLCM Haralick (5 props $\times$ 4 angles) & 20 & Rugosité, homogénéité, motifs \\
Forme & Moments de Hu ($\times$7) + métriques & 10 & Invariant rotation/translation/scale \\
\midrule
\textbf{Total} & & \textbf{566} & \\
\bottomrule
\end{tabular}
\caption{Composition du vecteur de features utilisé par les modèles ML}
\end{table}

Les features sont normalisées par \textbf{StandardScaler} ($\mu=0$, $\sigma=1$), 
fitté uniquement sur les données d'entraînement pour éviter la fuite de données.

% ============================================================
\section{Classification par Machine Learning}
% ============================================================

\subsection{Random Forest (RF)}

Ensemble de 300 arbres de décision avec bagging et sélection aléatoire de features.
\begin{itemize}
  \item \texttt{n\_estimators=300} : plateau de performance atteint (compromis précision/temps)
  \item \texttt{class\_weight='balanced'} : corrige le déséquilibre de classes (1\% à 14\%)
  \item \texttt{n\_jobs=-1} : parallélisation sur tous les cœurs CPU
\end{itemize}

\subsection{SVM à Noyau RBF}

Le SVM cherche l'hyperplan de marge maximale dans un espace de dimension infinie 
(noyau RBF : $K(x,y) = e^{-\gamma\|x-y\|^2}$).
\begin{itemize}
  \item \texttt{C=10} : régularisation standard pour features normalisées
  \item \texttt{gamma='scale'} : $\gamma = 1/(n\_\text{features} \times \text{Var}(X))$
  \item \texttt{probability=True} : Platt scaling pour les probabilités Top-3 (Streamlit)
\end{itemize}

\subsection{Résultats ML sur le Test Set}

\begin{table}[H]
\centering\small
\begin{tabular}{lcccccc}
\toprule
\textbf{Modèle} & \textbf{Accuracy} & \textbf{Précision} & \textbf{Rappel} & \textbf{F1} & \textbf{Train} & \textbf{Inférence} \\
\midrule
Random Forest & 89.04\% & 89.10\% & 89.04\% & 88.79\% & 5.5s & 0.09s \\
SVM (RBF)     & 87.33\% & 87.57\% & 87.33\% & 87.37\% & 44.6s & 5.4s \\
\bottomrule
\end{tabular}
\caption{Performances comparées RF vs SVM sur 2\,044 images de test}
\end{table}

\begin{figure}[H]
\centering
\begin{subfigure}{0.48\textwidth}
\fbox{\parbox{\textwidth}{\centering\vspace{2cm}
\textit{[confusion\_matrix\_rf.png]}
\vspace{2cm}}}
\caption{Matrice de confusion — Random Forest}
\end{subfigure}
\hfill
\begin{subfigure}{0.48\textwidth}
\fbox{\parbox{\textwidth}{\centering\vspace{2cm}
\textit{[confusion\_matrix\_svm.png]}
\vspace{2cm}}}
\caption{Matrice de confusion — SVM (RBF)}
\end{subfigure}
\caption{Matrices de confusion des modèles ML (12 classes, test set)}
\end{figure}

% ============================================================
\section{Extension Deep Learning}
% ============================================================

\subsection{Architecture MobileNetV2 (Transfer Learning)}

MobileNetV2 est un réseau convolutif léger pré-entraîné sur ImageNet (1,2M images, 
1000 classes). Nous adaptons ses 3,4M paramètres à notre problème agricole via 
le \textbf{Transfer Learning en 2 phases} :

\begin{enumerate}
  \item \textbf{Phase 1 — Feature Extraction} : couches MobileNetV2 gelées, 
  seule la tête Dense(12) est entraînée (\texttt{lr=1e-3}, 10 epochs).
  \item \textbf{Phase 2 — Fine-tuning} : les 30 dernières couches dégelées, 
  réentraînement avec \texttt{lr=1e-5} pour adaptation au domaine foliaire.
\end{enumerate}

La tête de classification ajoutée : GlobalAveragePooling2D $\rightarrow$ 
BatchNormalization $\rightarrow$ Dropout(0.3) $\rightarrow$ Dense(12, softmax).

\subsection{Résultats MobileNetV2}

\begin{table}[H]
\centering\small
\begin{tabular}{lccccc}
\toprule
\textbf{Modèle} & \textbf{Accuracy} & \textbf{Précision} & \textbf{Rappel} & \textbf{F1} & \textbf{Train} \\
\midrule
MobileNetV2 & \textbf{93.59\%} & \textbf{93.85\%} & \textbf{93.59\%} & \textbf{93.57\%} & $\sim$90 min \\
\bottomrule
\end{tabular}
\caption{Performances de MobileNetV2 sur 2\,044 images de test}
\end{table}

\begin{figure}[H]
\centering
\fbox{\parbox{0.9\textwidth}{\centering\vspace{1.5cm}
\textit{[training\_history.png — Courbes Accuracy \& Loss Phase 1 + Fine-tuning]}
\vspace{1.5cm}}}
\caption{Courbes d'apprentissage MobileNetV2}
\end{figure}

\subsection{Explicabilité — Grad-CAM}

L'algorithme Grad-CAM (Gradient-weighted Class Activation Mapping) visualise 
les régions de l'image qui ont le plus influencé la décision du réseau. 
Nous utilisons \texttt{tf.GradientTape} ciblant la couche \texttt{out\_relu} 
(sortie $7\times7\times1280$ de MobileNetV2).

\begin{figure}[H]
\centering
\fbox{\parbox{0.9\textwidth}{\centering\vspace{2cm}
\textit{[Captures Streamlit : Original | Heatmap pure | Overlay Grad-CAM]}
\vspace{2cm}}}
\caption{Visualisation Grad-CAM — zones décisives pour le diagnostic (rouge = haute activation)}
\end{figure}

% ============================================================
\section{Étude Comparative ML vs Deep Learning}
% ============================================================

\begin{table}[H]
\centering\small
\begin{tabular}{lcccccc}
\toprule
\textbf{Critère} & \textbf{Random Forest} & \textbf{SVM (RBF)} & \textbf{MobileNetV2} \\
\midrule
Accuracy          & 89.04\%   & 87.33\%   & \textbf{93.59\%} \\
F1-score          & 88.79\%   & 87.37\%   & \textbf{93.57\%} \\
Temps entraînement & 5.5s    & 44.6s     & $\sim$90 min \\
Inférence/image   & 0.04 ms  & 2.6 ms    & 17 ms \\
GPU requis        & Non       & Non       & Optionnel \\
Interprétabilité  & Features  & Opaque    & \textbf{Grad-CAM} \\
Mémoire modèle    & $\sim$7 Mo & $\sim$15 Mo & $\sim$14 Mo \\
\bottomrule
\end{tabular}
\caption{Benchmark comparatif des 3 modèles sur le même test set (2\,044 images)}
\end{table}

\begin{figure}[H]
\centering
\fbox{\parbox{0.9\textwidth}{\centering\vspace{1.5cm}
\textit{[benchmark\_chart.png — Graphe comparatif Accuracy \& F1-score]}
\vspace{1.5cm}}}
\caption{Visualisation du benchmark ML vs DL}
\end{figure}

\noindent\textbf{Analyse :} MobileNetV2 surpasse les modèles ML de \textbf{+4.55\%} 
en accuracy grâce à l'apprentissage automatique de features spatiales. 
Cependant, le Random Forest reste compétitif avec des ressources minimales 
(5.5s d'entraînement, 0.04ms/image), idéal pour les systèmes embarqués.

% ============================================================
\section{Extensions Avancées (Section IV)}
% ============================================================

\subsection{Interface Graphique — Application Streamlit}

Nous avons développé une interface web interactive complète (\texttt{app.py}) 
accessible via \texttt{streamlit run app.py} sur \texttt{localhost:8501}.

\textbf{Fonctionnalités :}
\begin{itemize}
  \item Sélection du modèle (RF / SVM / MobileNetV2) en temps réel
  \item Upload d'images et affichage des probabilités Top-3 avec barres colorées
  \item Visualisation Grad-CAM : heatmap pure + overlay côte à côte
  \item Segmentation HSV de la feuille isolée
  \item Conseils agronomiques adaptatifs selon la sévérité détectée
\end{itemize}

\begin{figure}[H]
\centering
\fbox{\parbox{0.9\textwidth}{\centering\vspace{2cm}
\textit{[Capture d'écran de l'interface Streamlit — résultat complet]}
\vspace{2cm}}}
\caption{Interface Streamlit — diagnostic de Tomato Septoria leaf spot}
\end{figure}

\subsection{Pipeline Avancé — Comparaison Classique vs Deep}

Notre approche hybride constitue elle-même une extension avancée :
\begin{itemize}
  \item \textbf{Classique} : Features manuelles (GLCM, Hu, HSV) + RF/SVM
  \item \textbf{Deep} : Features automatiques (conv features) + MobileNetV2
  \item \textbf{Explicabilité} : Grad-CAM offre une interprétabilité spatiale que les modèles ML classiques ne peuvent fournir
\end{itemize}

% ============================================================
\section{Conclusion}
% ============================================================

Nous avons conçu et implémenté un \textbf{système complet de détection automatique 
de maladies foliaires} sur 12 classes et 13\,563 images. Les résultats obtenus 
démontrent la supériorité du Deep Learning (MobileNetV2 : 93.59\%) sur le 
Machine Learning classique (RF : 89.04\%, SVM : 87.33\%), tout en soulignant 
les avantages pratiques des modèles ML en termes de vitesse et de ressources.

\textbf{Points forts de notre réalisation :}
\begin{enumerate}
  \item Pipeline complet et modulaire, documenté dans un Jupyter Notebook de 46 cellules
  \item 5 méthodes de segmentation implémentées et comparées
  \item 566-dimensional feature vector (RGB + HSV + GLCM + Hu Moments)
  \item Fine-tuning MobileNetV2 en 2 phases avec callbacks avancés
  \item Grad-CAM pour l'explicabilité, déployé dans l'interface Streamlit
  \item Code hébergé publiquement sur GitHub
\end{enumerate}

\textbf{Perspectives d'amélioration :}
\begin{itemize}
  \item Tester EfficientNetV2 ou ConvNeXt pour gain de précision
  \item Déploiement mobile via TFLite (quantification int8)
  \item Segmentation sémantique profonde (U-Net) pour localiser les lésions
  \item Extension à de nouvelles espèces végétales (poivron, raisin)
\end{itemize}

\vspace{1cm}
\noindent\textbf{Lien GitHub :} 
\url{https://github.com/bahaell/Maladies\_des\_plantes}

\vspace{1cm}
\noindent\rule{\linewidth}{0.4pt}\\[0.2cm]
\textit{Bahaeddine Ellouze \& Omar Chokri — Section IGL4 — Avril 2026}

\end{document}
